% Paper-level pipeline schematic, redesigned from first principles.
% Compile (from this directory):
%   xelatex -jobname=fig_pipeline_overview_tikz fig_pipeline_overview.tex
% Output PDF:  fig_pipeline_overview_tikz.pdf  (named to avoid overwriting
% the existing matplotlib fig_pipeline_overview.pdf).
%
% Design logic (HORIZONTAL-SPLIT variant):
%   PANEL A (left): manipulation -> shared architecture -> encoder-memory race
%     centred on h_2, the contested object.  Left half of figure.
%   PANEL B (right): frozen-h_2 bank read by 3 method families,
%     each with a finding chip.  Right half of figure, with the 3 method
%     panels stacked VERTICALLY and a vertical convergence brace.
\documentclass[tikz,border=4pt]{standalone}
\usepackage{mathptmx}
\usepackage{amsmath}
\usepackage{amssymb}
\usepackage{graphicx}
\usepackage{tikz}
\usetikzlibrary{positioning, arrows.meta, calc, fit, backgrounds, shapes.geometric, shapes.misc, decorations.pathreplacing, decorations.markings}

\begin{document}

% ---- palette ----
\definecolor{cBlind}{RGB}{75,75,75}
\definecolor{cCoarse}{RGB}{31,119,180}
\definecolor{cFovea}{RGB}{214,39,40}
\definecolor{cUni}{RGB}{44,160,44}
\definecolor{cInk}{RGB}{40,40,46}
\definecolor{cMute}{RGB}{120,120,128}
\definecolor{cBgVis}{RGB}{255,239,221}
\definecolor{cBgInt}{RGB}{222,237,247}
\definecolor{cBgH2}{RGB}{255,247,221}
\definecolor{cBgFreeze}{RGB}{245,243,237}
\definecolor{cProbe}{RGB}{31,119,180}
\definecolor{cGeom}{RGB}{148,103,189}
\definecolor{cInter}{RGB}{255,127,14}
\definecolor{cAccent}{RGB}{197,90,17}

\begin{tikzpicture}[
    >={Latex[length=2.4mm,width=2.0mm]},
    font=\rmfamily\small,
    every node/.style={inner sep=2pt, outer sep=0pt},
    box/.style={draw=cInk, line width=0.6pt, rounded corners=2pt, align=center},
    flow/.style={->, draw=cInk, line width=0.7pt},
    flowAcc/.style={->, draw=cAccent, line width=1.0pt},
    routeVis/.style={->, draw=cAccent, line width=1.4pt},
    routeInt/.style={->, draw=cCoarse, line width=1.4pt},
    lbl/.style={font=\rmfamily\scriptsize, text=cMute},
    chip/.style={draw=cMute, line width=0.4pt, rounded corners=1.2pt, fill=white, inner sep=2.0pt, font=\rmfamily\scriptsize},
    track/.style={draw=cMute, line width=0.4pt, dashed, rounded corners=2pt, fill=white},
]

% ============================================================
%  PANEL A (LEFT HALF). condition col(0) -> encoder(2.9) -> LSTM(5.6) -> h_2(7.9) -> pi(9.5)
% ============================================================

% ---- (A1) Conditions ladder ----
% Thumbnails are 26mm, spacing ~2.7 tikz units between row centres (partial revert: still tighter than original 3.3).
% Spec is split into INPUT (left of encoder, on input arrow) and OUTPUT (right of encoder, on
% output arrow); see (A2b)/(A5) below.  Blind row keeps a brief inline note since it bypasses.
\foreach \name/\file/\col/\y/\note in {%
    Blind/fig1_setup_blind/cBlind/5.0/{},%
    Coarse/fig1_setup_matched/cCoarse/2.3/{},%
    Foveated/fig1_setup_foveated_fix/cFovea/-0.4/{},%
    Uniform/fig1_setup_uniform/cUni/-3.1/{}%
}{%
    \node[draw=\col, line width=0.8pt, inner sep=0pt] (cond-\name) at (0, \y)
        {\includegraphics[width=26mm,height=26mm]{\file}};
    \node[font=\rmfamily\bfseries\small, text=\col, anchor=west]
        at ($(cond-\name.east)+(1.5mm, 4.0mm)$) {\name};
    \node[font=\rmfamily\itshape\fontsize{7pt}{8pt}\selectfont, text=cMute, anchor=west, align=left]
        at ($(cond-\name.east)+(1.5mm, -2.0mm)$) {\note};
}

\node[font=\rmfamily\bfseries\footnotesize, text=cInk, anchor=south west, align=left]
    at (-1.0, 6.5) {Visual sensor};

% Bandwidth axis (label \normalsize = 0.7x of previous \Large; placed OUTSIDE the arrow on
% the LEFT, rotated bottom->top in standard left-axis convention).
\draw[->, line width=0.7pt, draw=cMute]
    (-1.7, 6.3) -- (-1.7, -4.4);
\node[font=\rmfamily\normalsize\itshape, text=cMute, rotate=90, anchor=south]
    at (-2.0, 0.95) {encoder spatial bandwidth};

% ---- (A2) Encoder (TRAPEZOID, taller, centred among 3 sighted conditions) ----
% Centred at (5.0, -0.4); left edge spans 54mm vertically (matches new condition spacing),
% right edge 30mm vertically (matches LSTM input window); 18mm wide.
\coordinate (encTL) at (4.1, 2.3);
\coordinate (encBL) at (4.1, -3.1);
\coordinate (encTR) at (5.9, 1.1);
\coordinate (encBR) at (5.9, -1.9);
% drop shadow (offset 1mm down-right)
\filldraw[draw=none, fill=cInk!18]
    ($(encTL)+(0.6mm,-0.6mm)$) -- ($(encTR)+(0.6mm,-0.6mm)$)
    -- ($(encBR)+(0.6mm,-0.6mm)$) -- ($(encBL)+(0.6mm,-0.6mm)$) -- cycle;
\filldraw[draw=cInk, line width=0.6pt, fill=cBgVis]
    (encTL) -- (encTR) -- (encBR) -- (encBL) -- cycle;
\node[font=\rmfamily\small, text=cInk, align=center] (enc) at (5.0, -0.4)
    {ResNet-18\\encoder};
% define helper coordinates so subsequent code that expects (enc.west)/(enc.east) still works
\coordinate (enc-west) at (4.1, -0.4);
\coordinate (enc-east) at (5.9, -0.4);
\coordinate (enc-north) at (5.0, 2.0);

% --- (A2b) 3 separate colored input arrows: each condition -> encoder left edge ---
% Each arrow is STRICTLY HORIZONTAL: condition right edge -> encoder left edge at same y.
% Encoder is centred at y=-0.4 with left edge spanning [-3.1, 2.3], matching cond y-values exactly.
% Input spec sits italic 7pt in the condition's palette colour, above each arrow.
\coordinate (encInC) at (4.1, 2.3);  % Coarse entry (top of left edge)
\coordinate (encInF) at (4.1, -0.4); % Foveated entry (middle of left edge)
\coordinate (encInU) at (4.1, -3.1); % Uniform entry (bottom of left edge)
% Coarse (blue) - strict horizontal
\draw[->, draw=cCoarse, line width=0.9pt]
    (cond-Coarse.east) -- (encInC)
    node[midway, above, font=\rmfamily\itshape\fontsize{7pt}{8pt}\selectfont,
         text=cCoarse, inner sep=1pt] {$48^2$\,RGB}
    node[midway, below, font=\rmfamily\itshape\fontsize{6.5pt}{7.5pt}\selectfont,
         text=cCoarse, inner sep=1pt] {spatial collapse};
% Foveated (red) - strict horizontal
\draw[->, draw=cFovea, line width=0.9pt]
    (cond-Foveated.east) -- (encInF)
    node[midway, above, font=\rmfamily\itshape\fontsize{7pt}{8pt}\selectfont,
         text=cFovea, inner sep=1pt] {$256^2$\,blur}
    node[midway, below, font=\rmfamily\itshape\fontsize{6.5pt}{7.5pt}\selectfont,
         text=cFovea, inner sep=1pt] {peripheral blur};
% Uniform (green) - strict horizontal
\draw[->, draw=cUni, line width=0.9pt]
    (cond-Uniform.east) -- (encInU)
    node[midway, above, font=\rmfamily\itshape\fontsize{7pt}{8pt}\selectfont,
         text=cUni, inner sep=1pt] {$256^2$\,RGB}
    node[midway, below, font=\rmfamily\itshape\fontsize{6.5pt}{7.5pt}\selectfont,
         text=cUni, inner sep=1pt] {full spatial layout};

% --- (A2c) Blind bypass: dashed arrow ABOVE the encoder, direct to LSTM L0 input area ---
% (path defined later once lstm is placed; placeholder here is now removed.)

% ---- (A3) Non-visual sensors with icons ----
% Box positioned BELOW the LSTM container.
\node[box, fill=white, minimum width=30mm, minimum height=14mm]
    (sens) at (8.4, -3.6) {};
% Compute coordinates for 2x2 icon+label grid inside the box.
% Grid is shifted ~3.5mm left of box centre, and icon-to-text gap halved (3mm -> 1.5mm)
% so each icon hugs its label.
\coordinate (sCx) at (sens.center);
% Top row: GPS (left col), Compass (right col)
% Bottom row: Goal (left col), PrevAct (right col)
\coordinate (sTLi) at ($(sens.center)+(-10.5mm, 2.7mm)$); % top-left icon centre
\coordinate (sTLt) at ($(sens.center)+( -9.0mm, 2.7mm)$); % top-left text anchor (west)
\coordinate (sTRi) at ($(sens.center)+(  0.3mm, 2.7mm)$);
\coordinate (sTRt) at ($(sens.center)+(  1.8mm, 2.7mm)$);
\coordinate (sBLi) at ($(sens.center)+(-10.5mm,-2.7mm)$);
\coordinate (sBLt) at ($(sens.center)+( -9.0mm,-2.7mm)$);
\coordinate (sBRi) at ($(sens.center)+(  0.3mm,-2.7mm)$);
\coordinate (sBRt) at ($(sens.center)+(  1.8mm,-2.7mm)$);

% --- icon: GPS pin (Google-Maps style: rounded top + tapered point + inner dot) ---
% Path starts at the bottom point (0,-1.5mm), curves out to the rounded top
% circle (centred at y=+0.5mm, radius 0.9mm), wraps around it, and tapers back
% down to the bottom point.
\begin{scope}[shift={(sTLi)}, line cap=round, line join=round]
  \draw[draw=cInk, line width=0.4pt, fill=white]
       (0,-1.5mm)
       .. controls +(0.55mm, 0.55mm) and +(-0.5mm, -0.9mm) ..
       (0.9mm, 0.5mm)
       arc[start angle=0, end angle=180, radius=0.9mm]
       .. controls +(0.5mm, -0.9mm) and +(-0.55mm, 0.55mm) ..
       (0,-1.5mm) -- cycle;
  \fill[cInk] (0, 0.5mm) circle (0.35mm);
\end{scope}
\node[font=\rmfamily\scriptsize, text=cInk, anchor=west] at (sTLt) {GPS};

% --- icon: Compass (4-arm star inside circle) ---
\begin{scope}[shift={(sTRi)}, line cap=round]
  \draw[draw=cInk, line width=0.4pt] (0,0) circle (1.4mm);
  % 4 arms (diamond points)
  \fill[cInk] (0,1.4mm) -- (0.4mm,0) -- (0,-1.4mm) -- (-0.4mm,0) -- cycle;
  \fill[cInk] (1.4mm,0) -- (0,0.4mm) -- (-1.4mm,0) -- (0,-0.4mm) -- cycle;
\end{scope}
\node[font=\rmfamily\scriptsize, text=cInk, anchor=west] at (sTRt) {Compass};

% --- icon: Goal (3 concentric rings) ---
\begin{scope}[shift={(sBLi)}]
  \draw[draw=cInk, line width=0.35pt] (0,0) circle (1.4mm);
  \draw[draw=cInk, line width=0.35pt] (0,0) circle (0.9mm);
  \fill[cInk] (0,0) circle (0.35mm);
\end{scope}
\node[font=\rmfamily\scriptsize, text=cInk, anchor=west] at (sBLt) {Goal};

% --- icon: Action (curved arrow up-right) ---
\begin{scope}[shift={(sBRi)}, line cap=round]
  \draw[->, draw=cInk, line width=0.5pt, >={Latex[length=1.2mm,width=1.0mm]}]
       (-1.2mm,-1.0mm) .. controls +(0.6mm, 1.4mm) and +(-0.6mm, -0.6mm) ..
       (1.2mm, 1.2mm);
\end{scope}
\node[font=\rmfamily\scriptsize, text=cInk, anchor=west] at (sBRt) {PrevAct};

\node[lbl, anchor=south] at ($(sens.south)+(0, -3.6mm)$)
    {non-visual sensors};

% ---- (A4) LSTM (title split across 2 lines so it doesn't crowd L0/L1/L2) ----
% Centred at y=-0.4 so its left-edge centerline aligns with the encoder right-edge centerline,
% allowing all 3 colored encoder->LSTM arrows to be STRICTLY HORIZONTAL (no Manhattan bend).
\node[draw=cInk, line width=0.6pt, rounded corners=2pt, fill=white,
      minimum height=32mm, minimum width=20mm] (lstm) at (8.4, -0.4)
      {};
\node[font=\rmfamily\bfseries\small, text=cInk, anchor=north, align=center]
    at ($(lstm.north)+(0,-2mm)$) {3-layer\\LSTM};
\foreach \i/\lab in {0/L0, 1/L1, 2/L2}{
    \draw[draw=cMute, line width=0.3pt, fill=cBgInt!50, rounded corners=0.5pt]
        ($(lstm.south west)+(2.3mm, 2mm + \i*4.5mm)$) rectangle
        ($(lstm.south east)+(-2.3mm, 5.8mm + \i*4.5mm)$);
    \node[font=\rmfamily\scriptsize, text=cInk]
        at ($(lstm.south west)+(4.6mm, 3.9mm + \i*4.5mm)$) {\lab};
}

% ---- (A5) visual route: 3 separate colored output arrows from encoder to LSTM ----
% Each output arrow carries an italic spec label (NO feature-map glyphs).
% After LSTM was moved down to align with encoder right-edge centerline (y=-0.4), all 3
% encoder->LSTM arrows are STRICTLY HORIZONTAL (no Manhattan bend).
% Encoder right edge spans y in [-1.9, 1.1]; LSTM west spans y in [-2.0, 1.2] (encloses).
\coordinate (encOutC) at (5.9,  1.10); % Coarse exit (top of right edge)
\coordinate (encOutF) at (5.9, -0.40); % Foveated exit (mid)
\coordinate (encOutU) at (5.9, -1.90); % Uniform exit (bottom of right edge)
\coordinate (lstmInC) at (7.4,  1.10);
\coordinate (lstmInF) at (7.4, -0.40);
\coordinate (lstmInU) at (7.4, -1.90);

% --- Coarse output arrow: STRICT HORIZONTAL ---
\draw[->, draw=cCoarse, line width=0.9pt]
    (encOutC) -- (lstmInC);
\node[font=\rmfamily\itshape\fontsize{7pt}{8pt}\selectfont, text=cCoarse,
      anchor=south, inner sep=1pt]
    at (6.65, 1.18) {$\to 1{\times}1$};
% --- Foveated output arrow: STRICT HORIZONTAL ---
\draw[->, draw=cFovea, line width=0.9pt]
    (encOutF) -- (lstmInF);
\node[font=\rmfamily\itshape\fontsize{7pt}{8pt}\selectfont, text=cFovea,
      anchor=south, inner sep=1pt]
    at (6.65, -0.32) {$\to 8{\times}8$};
% --- Uniform output arrow: STRICT HORIZONTAL ---
\draw[->, draw=cUni, line width=0.9pt]
    (encOutU) -- (lstmInU);
\node[font=\rmfamily\itshape\fontsize{7pt}{8pt}\selectfont, text=cUni,
      anchor=south, inner sep=1pt]
    at (6.65, -1.82) {$\to 8{\times}8$};

% "visual route" label sits ABOVE the encoder top
\node[font=\rmfamily\bfseries\scriptsize, text=cAccent, anchor=south]
    at (5.0, 2.45) {visual route};

% Integration route: sens (below LSTM) -> UP into LSTM bottom (Manhattan: strict vertical)
\coordinate (sensOut) at (sens.north);
\draw[routeInt, rounded corners=1pt] (sensOut) -- (lstm.south);
\node[font=\rmfamily\bfseries\scriptsize, text=cCoarse, anchor=east]
    at ($(sens.north)+(-1mm, 4mm)$) {integration route};

% ---- (A5b) Blind bypass: dashed direct arrow from Blind thumbnail -> LSTM top area,
% routed ABOVE the encoder so it visually skips it. Strict Manhattan: H -> V -> H.
\draw[->, draw=cBlind, line width=0.8pt, dashed, rounded corners=1pt]
    (cond-Blind.east) -| ($(lstm.west)+(-4mm, 15.5mm)$) -- ($(lstm.west)+(0, 15.5mm)$);
\node[font=\rmfamily\scriptsize\itshape, text=cBlind, anchor=south, inner sep=1pt]
    at (5.0, 5.05) {bypass};
% Blind role label: italic 6.5pt, in palette colour, just below the bypass
% horizontal segment near the Blind thumbnail (the bypass runs at y=5.0).
\node[font=\rmfamily\itshape\fontsize{6.5pt}{7.5pt}\selectfont, text=cBlind,
      anchor=north west, inner sep=1pt]
    at ($(cond-Blind.east)+(1.5mm, -1.0mm)$) {non-visual baseline};

% ---- (A6) race caption removed ----

% ---- (A7) h_2 hub ----
\node[draw=cAccent, line width=1.2pt, fill=cBgH2, rounded corners=2pt,
      minimum width=16mm, minimum height=20mm] (h2) at (10.7, -0.4) {};
\node[font=\rmfamily\bfseries\Large, text=cAccent] at (h2.center) {$\mathbf{h}_2$};

\draw[flowAcc, line width=1.2pt] (lstm.east) -- (h2.west);

% ---- (A8) policy head (positioned BELOW h2, vertically aligned) ----
\node[draw=cInk, line width=0.6pt, fill=white, regular polygon, regular polygon sides=3,
      shape border rotate=-90, minimum size=8mm] (pi) at (10.7, -2.4) {};
\node[font=\rmfamily\bfseries] at (pi.center) {$\pi$};
% h2 -> pi: strict vertical Manhattan from h2 south to pi north
\draw[flow] (h2.south) -- (pi.north);
% Action arrow: stub DOWN from pi south, then RIGHT to the Behavioral panel (T3) left edge.
% Stub continues down well below T3's mid-height so the |- routing is clearly DOWN-then-RIGHT.
\coordinate (actStub) at ($(pi.south)+(0, -8mm)$);
\draw[flow, line width=0.7pt] (pi.south) -- (actStub)
    node[midway, right, font=\rmfamily\scriptsize, text=cInk, inner sep=1pt] {action};

% ---- (A9) Train/eval footnote removed ----

% (A) section label  -- top-left of A.  0.7x of previous \LARGE -> \large.
\node[font=\rmfamily\bfseries\large, text=cInk, anchor=west]
    at (-1.7, 7.6) {(A)\enspace Training Architecture};

% ============================================================
%  PANEL B (RIGHT HALF) - 3 method panels stacked vertically
%  Right half occupies x in [13.5, 21.5]
% ============================================================

% vertical divider between A and B (shortened bottom to match content extent)
\draw[draw=cMute, line width=0.3pt, dotted]
    (13.8, 7.4) -- (13.8, -4.0);

% (B) section label -- top-left of B.  0.7x of previous \LARGE -> \large.
\node[font=\rmfamily\bfseries\large, text=cAccent, anchor=west]
    at (14.1, 7.2) {(B)\enspace Memory Analysis};

% ---- (B1) Bank (narrower, shifted DOWN/right; visual hint = 4 stacked condition stripes) ----
% Width: 36mm (~62% of previous 58mm).  Centred at (17.6, 6.0) -- moved down from 6.6 to give
% breathing room from the (B) label and slightly right.
% Layout: top row holds the bold title "h_2 bank" and the 4 coloured stripes (one per
% condition, hint of stored vectors); bottom row holds the italic "500 ep x 4 cond" subtitle.
\node[draw=cAccent, line width=0.8pt, fill=cBgFreeze, rounded corners=2pt,
      minimum width=36mm, minimum height=12mm] (bank) at (16.95, 5.6) {};
\node[font=\rmfamily\bfseries\small, text=cAccent, anchor=center]
    at ($(bank.center)+(0, 2.0mm)$) {$\mathbf{h}_2$ bank};
\node[font=\rmfamily\scriptsize\itshape, text=cMute, anchor=center, align=center]
    at ($(bank.center)+(0, -2.0mm)$)
    {500 ep $\times$ 4 cond.};

% extract arrow: h2 -> right -> up -> bank (Manhattan: V up then H right)
\draw[->, draw=cAccent, line width=0.9pt, rounded corners=1pt]
    (h2.north) |- (bank.west);
\node[font=\rmfamily\bfseries\scriptsize, text=cAccent, anchor=south west, align=left]
    at ($(h2.north east)+(0.5mm, 0.5mm)$) {freeze\\\& extract};

% ---- (B3) three method tracks - stacked VERTICALLY in right half ----
% Each panel: ~70mm wide, ~22mm tall.  Centred at x=15.
% Y positions: T1 top, T2 middle, T3 bottom

% Track 1 (TOP): linear probes
\node[track, fill=cProbe!8, minimum width=58mm, minimum height=22mm]
    (T1) at (16.95, 3.1) {};
\node[anchor=north west, align=left, inner sep=0pt]
    at ($(T1.north west)+(2mm,-1.5mm)$)
    {{\rmfamily\bfseries\normalsize\color{cProbe} Decodability}\\[-0.4mm]{\rmfamily\itshape\fontsize{9pt}{10pt}\selectfont\color{cProbe!75!white} (probes)}};
\node[font=\rmfamily\fontsize{7pt}{8pt}\selectfont\itshape, text=cInk, anchor=north west,
      align=left, text width=50mm, inner sep=0pt]
    at ($(T1.north west)+(2mm,-10mm)$) {``what does $\mathbf{h}_2$ encode?''};
% Chip on the bottom-left of T1
\node[chip, text=cProbe, font=\rmfamily\bfseries\scriptsize, align=left, anchor=south west]
    at ($(T1.south west)+(2mm, 1.3mm)$) {\textbf{H1}\,: bandwidth--allocation tradeoff};

% Track 2 (MIDDLE): geometry
\node[track, fill=cGeom!8, minimum width=58mm, minimum height=22mm]
    (T2) at (16.95, 0.25) {};
\node[anchor=north west, align=left, inner sep=0pt, text width=58mm]
    at ($(T2.north west)+(2mm,-1.5mm)$)
    {{\rmfamily\bfseries\normalsize\color{cGeom} Intrinsic Structure}\\[-0.4mm]{\rmfamily\itshape\fontsize{9pt}{10pt}\selectfont\color{cGeom!75!white} (representational geometry)}};
\node[font=\rmfamily\fontsize{7pt}{8pt}\selectfont\itshape, text=cInk, anchor=north west,
      align=left, text width=50mm, inner sep=0pt]
    at ($(T2.north west)+(2mm,-10mm)$) {``do conditions share a code?''};
\node[chip, text=cGeom, font=\rmfamily\bfseries\scriptsize, align=left, anchor=south west]
    at ($(T2.south west)+(2mm, 1.3mm)$) {\textbf{H2}\,: non-interchangeable subspaces};

% Track 3 (BOTTOM): intervention
\node[track, fill=cInter!8, minimum width=58mm, minimum height=22mm]
    (T3) at (16.95, -2.6) {};
\node[anchor=north west, align=left, inner sep=0pt]
    at ($(T3.north west)+(2mm,-1.5mm)$)
    {{\rmfamily\bfseries\normalsize\color{cInter} Causal Intervention}\\[-0.4mm]{\rmfamily\itshape\fontsize{9pt}{10pt}\selectfont\color{cInter!75!white} (behavior)}};
\node[font=\rmfamily\fontsize{7pt}{8pt}\selectfont\itshape, text=cInk, anchor=north west,
      align=left, text width=50mm, inner sep=0pt]
    at ($(T3.north west)+(2mm,-10mm)$) {``does the policy rely on $\mathbf{h}_2$?''};
\node[chip, text=cInter, font=\rmfamily\bfseries\scriptsize, align=left, anchor=south west]
    at ($(T3.south west)+(2mm, 1.3mm)$) {availability $\neq$ consumption};

% bank -> tracks arrows (vertical, all going down through stack)
\draw[->, draw=cMute, line width=0.4pt] (bank.south) -- ($(T1.north)+(0, 0)$);
\draw[->, draw=cMute, line width=0.4pt] (T1.south) -- ($(T2.north)+(0, 0)$);
\draw[->, draw=cMute, line width=0.4pt] (T2.south) -- ($(T3.north)+(0, 0)$);

% (Removed: convergence brace + "3 methods => same contrast" text - the 3 stacked
% panels make convergence self-evident.)

% ---- Action -> Behavioral (Intervention) panel ---- (Manhattan: DOWN then RIGHT to T3 west) ----
% From actStub (well below pi), arrow runs RIGHT to T3 west bottom.
\coordinate (behavEntry) at ($(T3.west)+(0, -6mm)$);
\draw[->, draw=cInter, line width=0.9pt, dashed, rounded corners=1pt]
    (actStub) |- (behavEntry);
% "behavioral" label on the horizontal segment, well to the right of the action stub
\node[font=\rmfamily\scriptsize\itshape, text=cInter, inner sep=1pt, anchor=south]
    at (12.6, -3.2) {behavioral};

\end{tikzpicture}
\end{document}
